La congestión del tráfico urbano representa uno de los problemas más persistentes y
costosos de las ciudades contemporáneas. El crecimiento sostenido del parque automotor,
combinado con una infraestructura vial que raramente crece al mismo ritmo, genera demoras
acumuladas que impactan directamente en la productividad económica, la calidad del aire y
la calidad de vida de los habitantes \cite{papageorgiou2003review,garber2019traffic}. En
ciudades de tamaño medio de América Latina, donde los recursos para grandes obras de
infraestructura son escasos, la optimización del control semafórico existente se presenta
como una de las intervenciones de mayor relación costo-beneficio para mejorar la fluidez
del tráfico sin requerir expansión física de la red vial \cite{stevanovic2010adaptive}.

Los sistemas de control semafórico tradicionales basados en tiempos fijos o predeterminados
presentan limitaciones fundamentales frente a la variabilidad inherente de la demanda urbana.
Un plan de señales diseñado para el pico de la mañana es ineficiente en el valle del mediodía;
un plan calibrado para días laborables falla en eventos o festividades. Esta rigidez
operativa motiva el desarrollo de estrategias adaptativas capaces de ajustar los parámetros
de señalización (duración de ciclo, distribución de tiempos verdes, desplazamientos entre
intersecciones) en respuesta a las condiciones observadas en tiempo real
\cite{papageorgiou2003review,stevanovic2010adaptive,li2016realtime}.

La convergencia de tres tendencias tecnológicas abre nuevas posibilidades para abordar este
problema. En primer lugar, la disponibilidad masiva de sensores de bajo costo y plataformas
de simulación de tráfico de código abierto, como SUMO (Simulation of Urban MObility),
permite evaluar y validar estrategias de control bajo condiciones controladas y reproducibles
antes de su despliegue en campo \cite{behrisch2011sumo}. En segundo lugar, las técnicas de
inteligencia computacional, en particular la lógica difusa y los algoritmos de optimización
metaheurística, ofrecen mecanismos para modelar el conocimiento experto sobre congestión y
para buscar configuraciones semafóricas óptimas en espacios de decisión complejos, sin requerir
modelos analíticos exactos de la demanda \cite{zadeh1965fuzzy,kennedy1995particle}. En tercer
lugar, las tecnologías de almacenamiento descentralizado (IPFS y redes BlockDAG) permiten
registrar los datos generados por el sistema de forma inmutable y auditable, sin depender de
un operador central de confianza, lo que resulta especialmente relevante para infraestructuras
de gestión pública donde la trazabilidad y la rendición de cuentas son requisitos
institucionales \cite{benet2014ipfs,sompolinsky2021phantom}.

\section{Motivación y contexto}

En el contexto paraguayo, la gestión del tráfico urbano en ciudades intermedias y en el área
metropolitana de Asunción enfrenta el desafío adicional de operar con presupuestos municipales
limitados y con infraestructura tecnológica heterogénea. Las soluciones comerciales de control
adaptativo, ampliamente desplegadas en ciudades de países desarrollados, suelen implicar costos
de licenciamiento, dependencia de proveedores y dificultades de integración con sistemas
legados, barreras que frenan su adopción a escala local. Este escenario plantea la necesidad
de explorar alternativas basadas en software libre, componentes de hardware de consumo y
arquitecturas abiertas que puedan desplegarse y mantenerse con los recursos técnicos disponibles
en municipios de la región.

Paralelamente, la creciente adopción de paradigmas de smart city en la planificación
urbana de América Latina genera demanda de soluciones que no solo optimicen parámetros
operativos, sino que también garanticen la integridad y disponibilidad de los datos generados.
En sistemas de gestión pública, la capacidad de auditar retrospectivamente el comportamiento
del sistema (verificar qué configuraciones se aplicaron, cuándo y con qué efecto) constituye
un requisito de transparencia que las plataformas centralizadas tradicionales satisfacen de
forma insuficiente.

\section{Objetivos}

\subsection{Objetivo general}

Diseñar, implementar y evaluar experimentalmente un sistema modular y descentralizado para la
optimización adaptativa del control semafórico urbano y la persistencia verificable de los
datos asociados, utilizando técnicas de inteligencia computacional y tecnologías de
almacenamiento distribuido sobre hardware de consumo estándar.

\subsection{Objetivos específicos}

\begin{itemize}
    \item Desarrollar un módulo de simulación de tráfico que genere escenarios controlados
          y reproducibles mediante SUMO, incluyendo detección automática de cuellos de botella
          y exposición de métricas agregadas por intersección semafórica.

    \item Diseñar e implementar un módulo de análisis y optimización que clasifique el estado
          de congestión mediante inferencia difusa tipo Mamdani y ajuste los tiempos semafóricos
          mediante optimización por enjambre de partículas (PSO).

    \item Construir un módulo de persistencia verificable que almacene datos de tráfico y
          resultados de optimización en IPFS, registrando los identificadores de contenido en
          una red BlockDAG mediante contratos inteligentes para garantizar inmutabilidad y
          trazabilidad.

    \item Integrar los módulos anteriores mediante un servicio de orquestación con interfaz
          REST que coordine el flujo completo desde la recepción de observaciones hasta la
          confirmación del almacenamiento distribuido.

    \item Evaluar experimentalmente el desempeño del sistema en términos de escalabilidad,
          consumo de recursos, latencia de almacenamiento y costo operativo, comparando con
          alternativas basadas en Ethereum y plataformas centralizadas IoT--cloud.
\end{itemize}

\section{Contribuciones}

El presente trabajo realiza las siguientes contribuciones:

\begin{enumerate}
    \item Una arquitectura de microservicios de código abierto que integra simulación de
          tráfico, optimización difusa-PSO y almacenamiento descentralizado en un flujo
          operativo unificado, demostrada en hardware de consumo estándar.

    \item Un protocolo de persistencia verificable para datos de tráfico urbano basado en
          IPFS y contratos inteligentes sobre BlockDAG, con confirmación en $\approx$1~s y
          costo inferior a \$1{,}2$\times$10\textsuperscript{-5}~USD por transacción,
          superando en más de cuatro órdenes de magnitud el costo operativo de Ethereum.

    \item Evidencia cuantitativa de la viabilidad del sistema para escenarios de hasta 2000
          sensores en hardware embebido de bajo costo, con tiempo de procesamiento
          $\approx 0{,}014$~s por sensor y consumo de memoria estable entre 198 y 219~MB.

    \item Un conjunto de experimentos reproducibles, con configuraciones, datos y métricas
          documentados, que sirven de línea base para investigaciones futuras en optimización
          de tráfico con requisitos de trazabilidad verificable.
\end{enumerate}

Los hallazgos principales fueron publicados en el artículo \emph{Scalable and Decentralized
Urban Traffic Optimization with Verifiable Storage for Smart City Deployments}, presentado en
la 2025 IEEE CHILEAN Conference on Electrical, Electronics Engineering, Information and
Communication Technologies (CHILECON~2025), Valparaíso, Chile \cite{chilecon2025}.

\section{Estructura del documento}

El resto del documento se organiza de la siguiente manera. El Capítulo~2 presenta el marco
teórico que fundamenta el trabajo: ingeniería de tráfico y control semafórico, lógica difusa
y sistemas de inferencia Mamdani, optimización por enjambre de partículas, IPFS como sistema
de almacenamiento distribuido por contenido, y las estructuras BlockDAG y blockchain
empleadas para el registro inmutable. El Capítulo~3 describe la visión general del sistema,
su arquitectura modular y el flujo de información entre componentes. Los Capítulos~4 a 7
detallan el diseño e implementación de cada módulo: simulación traffic-sim,
optimización traffic-sync, almacenamiento verificable traffic-storage y
orquestación traffic-control. El Capítulo~8 presenta la evaluación experimental con
los resultados de escalabilidad, latencia y costo. Finalmente, el Capítulo~9 recoge las
conclusiones, limitaciones y líneas de trabajo futuro.
